\section{Future Work: Possible Improvements}
\label{sec:future}

The improvements proposed in this section correspond either to a limitation listed in Section~\ref{subsec:limitations} or to one of the measured results. They are ranked in Table~\ref{tab:roadmap} according to the benefit we expect for the effort involved.

\subsection{The Optimization Study}
\label{subsec:fw-opt}

With three training seeds per arm, the intervals would also include training noise, and with several search seeds per method it would become clear whether the PSO result is repeatable. The small proxy differences among the top settings disappeared after full training. A multi-fidelity search such as successive halving or Hyperband~\cite{li2018hyperband} could handle this by promoting the best candidates to longer runs before a winner is selected. Since the swarm converged early, and in \gentwo did not improve at all after round~1, PSO variants that maintain exploration are worth trying. Examples are the constriction factor of Clerc and Kennedy~\cite{clerc2002particle}, a ring topology, an adaptive inertia weight, or simply more particles and rounds.

Other optimizers can be added as further arms against the same baselines with 32 evaluations, e.g., a genetic or evolutionary algorithm~\cite{young2015optimizing}, differential evolution~\cite{storn1997differential} and Bayesian optimization such as the tree-structured Parzen estimator of Optuna~\cite{akiba2019optuna}. This also answers the question of why PSO was chosen and not, for instance, the beluga whale optimizer of Sayed et al.~\cite{sayed2024intelligent}. The protocol does not depend on the optimizer, and its value lies in confronting every candidate with the same baselines. A larger search space that includes color augmentation, translation, loss gains and image size would show whether six settings were too few to beat the defaults. A domain-aware fitness computed on real-world validation images could address the fact that the tuning gains did not reach real-world photographs. The proxy check should also be repeated for \gentwo with a 100-epoch reference and a run of the defaults, which would test the hypothesis of Section~\ref{subsec:why}. On the statistical side, Holm's correction, intervals for per-class AP and stored bootstrap resamples would strengthen the analysis.

\subsection{Data and the Indian Domain}
\label{subsec:fw-data}

The most important next step is to obtain Indian training data with per-item boxes, either by relabeling a part of DWSD or through a new collection with a held-out test portion. This would address the failure on DWSD as well as the labeling mismatch that is partly responsible for it. Because the real-world data of \gentwo reduced the studio accuracy (Section~\ref{subsec:deployment}), new data should be added through domain-aware training, for example by retaining all studio images or by using a larger model. The weak OTHER class requires more examples, or it should be divided into clearer classes or merged with existing ones. Clutter-oriented augmentations such as copy-paste~\cite{ghiasi2021simple}, smaller object scales, blur and lighting changes are aimed at the small, overlapping and partly hidden items seen in real-world photographs. The classes could also be aligned with the three streams required by the SWM Rules 2016, namely bio-degradable, non-bio-degradable and domestic hazardous~\cite{moefcc2016solid}, so that every detection corresponds directly to a bin.

\subsection{The Model}
\label{subsec:fw-model}

Training larger or newer detectors (YOLOv8s, YOLOv8m, YOLO11~\cite{jocher2024yolo11}) with the same settings would reveal how much of the real-world gap is caused by limited model capacity, though at the expense of speed on edge devices. Sliced inference on overlapping tiles~\cite{akyon2022slicing} is suited to wide photographs with many small items. Confidence calibration, e.g., temperature scaling~\cite{guo2017calibration}, would make a score of 0.5 correspond to roughly 50\% correct detections, which is important when a bin acts on a threshold.

\subsection{Combining the Three Models}
\label{subsec:fw-ensemble}

According to the point estimates, no single arm is the best on every \gentwo test set (Table~\ref{tab:gen2-test}). On real-world photographs the three \mapfifty values are within 0.012 of each other, but the arms differ more in the details. Precision at the maximum-F1 confidence varies from 0.366 (B) to 0.556 (C), the AP@50 of PAPER varies from 0.239 (A) to 0.330 (C), and random search is the best on GLASS (0.296). These values have no intervals and the agreement between the models was not measured. Still, they suggest that the models make partly different errors, and a combination of the three arms, which requires no additional training, might perform better than any one of them.

Three ways of combining them can be compared. The first is box fusion, in which the boxes of A, B and C are pooled and the overlapping ones are merged by weighted boxes fusion~\cite{solovyev2021weighted} or non-maximum suppression. The second is majority voting, where a box is kept only when at least two models find it (IoU $\ge 0.5$). This should raise precision, and boxes found by a single model can be flagged as uncertain rather than thrown away. The third is weighted voting according to the validation \mapfifty of each model. Test scores must not be used as weights, as that would leak the test split. From the soft-computing point of view, a search produces a whole population of good settings (the final swarm and the best random samples). Training the top-$k$ settings and combining them would turn the swarm into an ensemble generator. The converged swarm may not be diverse enough, so the best random samples should be mixed in. Every ensemble should be compared with the best single arm, selected on validation data, using the same paired bootstrap on each test set. The price is inference time, which becomes about three times longer: roughly 100~ms per image on the laptop CPU, as against 31 to 38~ms for a single \genone model.

\subsection{Deployment}
\label{subsec:fw-deploy}

The single confidence threshold could be replaced by per-class thresholds chosen on validation data, because real-world precision is low and AP differs widely between classes. Our latency figures come only from a laptop CPU and a cloud GPU. The models should therefore be exported to ONNX or TFLite and timed on a Raspberry Pi, which is the kind of hardware used by Alshammari et al.~\cite{alshammari2025smart} for their smart-bin classifier. As actual sorting is a continuous process, the models should also be tested on live video.

\begin{table}[!t]
\centering
\caption{Future-Work Roadmap, Ordered by Expected Value for Effort}
\label{tab:roadmap}
\footnotesize
\setlength{\tabcolsep}{3pt}
\renewcommand{\arraystretch}{1.1}
\begin{tabular}{@{}c>{\raggedright\arraybackslash}p{2.75cm}>{\raggedright\arraybackslash}p{3.05cm}>{\raggedright\arraybackslash}p{1.45cm}@{}}
\toprule
Priority & Improvement & Addresses & Effort \\
\midrule
1 & Indian training data with per-item boxes & DWSD failure (0.023--0.027 \mapfifty) & Large \\
2 & Ensemble of A, B and C (fusion or voting) & Arms differ per class and test set; no arm wins everywhere & Small (no training) \\
3 & Three training seeds per arm, several search seeds & Tuning effect vs.\ training noise & Medium \\
4 & Domain-aware fitness (real-world validation) & Tuning gains did not carry to real-world photos & Medium \\
5 & Multi-fidelity search (Hyperband) & Proxy gaps vanish at 100 epochs & Medium \\
6 & GA, DE and Bayesian optimization at the same budget & ``Why PSO?'' & Medium \\
7 & Fix the OTHER class & AP@50 0.050--0.063 & Medium \\
8 & Larger detector; sliced inference & Small items, look-alike materials & Medium \\
9 & Edge export and latency on a Raspberry Pi & Real smart-bin hardware & Medium \\
10 & Swarm as an ensemble generator (top-$k$ settings) & Uses the whole PSO population & Medium \\
\bottomrule
\end{tabular}
\\[2pt]
\parbox{\columnwidth}{\footnotesize Effort: Small = hours, no GPU training; Medium = about a day with some GPU time; Large = several days or a new data pipeline. GA = genetic algorithm; DE = differential evolution.}
\end{table}
